\paragraph{How much of the vacancy stock can applicants reach?}
We applied Theorem~\ref{thm:residualreferral} after the rank-informed offer network's counterfactual allocation of five programs in the 2026 Chilean university cohort. It left 30,866 applicants without a seat and 21,061 seats vacant. Only 160 of these applicants had submitted any eligible program with an unfilled seat. Within that group, 24 shared just 11 reachable seats. This set witnesses the entire deficit of 13. A maximum flow therefore placed 147 additional people while protecting the provisional recipients. Their median new program was fourteenth on their submitted lists. A follow-up that both maximized cardinality and minimized the sum of entrants' reported ranks moved that median to thirteenth. If referrals must appear among an applicant's top five, ten, or fifteen submitted programs, the exact maxima fall to 4, 61, and 109 respectively. Willingness to accept a lower choice was not observed. Of the 147, seven had been unmatched even under full-list DA and three lacked both a full-list placement and an official regular offer. The other gains chiefly repaired losses created by the original shortlist.

Suppose instead each of the 160 applicants added their favorite program with a vacancy. This placed 156 previously unmatched people, displaced 11 tentative recipients and left 15 tentative recipients with a worse assignment, counting placement loss as worst. In contrast, serially reserving seats placed 144 with zero incumbent losses. DA run as a separate submarket among the eligible applicants and vacancies placed 143 with zero losses. We use previously submitted choices there, unlike the inferred additional schools in \citet{lepesure}. The maximum extension places four more than residual-submarket DA; the median new program ranks were fourteen and thirteen respectively. Maximizing the number of new placements need not select entrants' preferred alternatives.

The first-stage rule makes a large difference to what remains. After the first five submitted choices, the same construction added 5,691, 6,641 and 7,268 placements in the 2024, 2025 and 2026 cohorts. After the 2026 rank-blind network it added only 61. A separate optimistic bound allows incumbents to switch to a weakly preferred program while ignoring the stability constraints on their replacements. It allows at most 204 entrants after the informed network and 2,305 after the blind network, compared with 147 and 61 when all original contracts are held in the witness. Computing this bound starts from the incumbent matching and augments along paths retaining every matched applicant and every incumbent's original ordinal floor. Both the attained flow and its minimum cut were checked.

The rule requires the actual provisional assignment. After replacing priorities by noisy forecasts (independent score errors with standard deviations of 10 and 30 PAES points), choosing apparently vacant slots and then clearing with original scores displaced five and four provisional recipients. These are stress tests with artificial score errors. They do not estimate the accuracy of actual admissions forecasts.

\paragraph{Who can benefit from a second recommendation?}
After the contextual top-five rule in the 2,304 simulated NYC markets, reserved referrals produced 9,749 additional placements with no loss of any tentative recipient. Independent favorite-vacancy referrals produced 5,761; residual-market DA gave 8,460. The allocation among new recipients is still a policy decision. Among maximum-size extensions, favoring lower-income renters increased the number of new lower-income placements from 3,059 to 3,155 at the same total. Under this maximum the ordinary summer and winter panels produced only 13 and 51 additional lower-income placements, against 1,451 and 1,705 additional placements overall. In the deliberately enriched panels, the corresponding numbers were 1,431 of 3,127 and 1,660 of 3,466. Lower-income targeting redistributes the gains that the feasible rental graph permits. All figures are simulated renter instances; the two sampling designs describe different populations.
